% Shared mathematical content for the printable and searchable blueprints.
% Statements and mathematical proofs follow the source paper; Lean links and
% formalization maps record the verified implementation.

\begin{abstract}
Definitions, statements, and proofs follow the source paper in order.
Each formalized item links to its verified Lean declaration; the compact
formalization map beneath it identifies the implementation file.
\end{abstract}

\setcounter{section}{1}
\section{Functions and multiplicities}

Throughout, $n\in\N$ means $n\geq1$.  The max-plus operations are
\[
  a\oplusT b=\max\{a,b\},\qquad
  a\otimesT b=a+b,\qquad
  a\oslashT b=a-b.
\]
\begin{definition}[Tropicalization]
\paperlabel{def:nth-tropicalization}
\lean{NthTropicalNevanlinna.NthTropicalizesTo}
\leanok
For a positive family $q(X,\varepsilon)$ and $n\geq1$, write
\[
 q\xrightarrow[n]{\mathrm{trop}}Q
 \quad\Longleftrightarrow\quad
 \lim_{\varepsilon\to0^+}\varepsilon^n\log q(X,\varepsilon)=Q(X)
 \quad\text{for every }X\in\R.
\]
For finite nonempty supports $I,J$, degrees $u_i,v_j$, and real
coefficients $A_i,B_j$, the dequantized rational factor is
\[
 R_\varepsilon(X)=
 \frac{\sum_{i\in I}\exp((u_iX+A_i)/\varepsilon)}
      {\sum_{j\in J}\exp((v_jX+B_j)/\varepsilon)}.
\]
Its first tropicalization is
\[
 R(X)=\max_{i\in I}(u_iX+A_i)-\max_{j\in J}(v_jX+B_j).
\]
An active product of $n$ such factors, with the source parameter scaled so
that the logarithm is multiplied by $\varepsilon^n$, is the paper's
$n$-th tropical source.
\end{definition}

\begin{proposition}[Formalized tropicalization rules]
\paperlabel{prop:tropicalization-rules}
\lean{NthTropicalNevanlinna.activeDequantizedNthTropicalSource_tropicalizesTo,NthTropicalNevanlinna.nthTropicalizesTo_add,NthTropicalNevanlinna.nthTropicalizesTo_mul,NthTropicalNevanlinna.nthTropicalizesTo_div}
\leanok
\uses{def:nth-tropicalization}
Finite log-sum-exp tropicalizes to maximum.  Rational numerator/denominator
tropicalizes to the difference of their maxima, an active product of $n$
factors tropicalizes to the sum of the factor limits, and tropicalizable
positive families satisfy
\[
 f+g\mapsto F\oplusT G,\qquad fg\mapsto F\otimesT G,
 \qquad f/g\mapsto F\oslashT G.
\]
\end{proposition}

\begin{proof}
\leanok
If $M=\max_i a_i$, then
\[
 M\leq\varepsilon\log\sum_i e^{a_i/\varepsilon}
 \leq M+\varepsilon\log|I|,
\]
so the squeeze theorem gives the finite log-sum-exp limit.  Apply this once
to the numerator and once to the denominator and subtract.  Logarithms turn a
finite product into a finite sum, so the $n$ active factor limits add.  For
addition, factor out the larger positive summand; the remaining correction is
between $0$ and $\varepsilon^n\log2$ and therefore vanishes.  The product and
quotient identities follow directly from $\log(fg)=\log f+\log g$ and
$\log(f/g)=\log f-\log g$.
\end{proof}

\begin{definition}[$n$-th tropical meromorphic functions]
\paperlabel{def:nth-meromorphic}
\lean{NthTropicalNevanlinna.PolynomialPresentation,NthTropicalNevanlinna.IsNthTropicalMeromorphicFunction,NthTropicalNevanlinna.NthTropicalMeromorphicFunction}
\leanok
An \emph{$n$-th tropical meromorphic function} is a continuous function
$f:\R\to\R$ admitting a locally finite, bi-infinite increasing mesh
$(x_i)_{i\in\mathbb Z}$ and polynomial pieces $p_i$ such that
\[
 f(x)=p_i(x)\quad(x\in[x_{i-1},x_i]),\qquad
 \deg p_i\leq n,
\]
and at least one piece has degree $n$.  The mesh is auxiliary; its redundant
cut points are not singularities.
\end{definition}

\begin{leanproof}
\leanfile{PolynomialPresentation}, \leanfile{IsNthTropicalMeromorphicFunction}, and
\leanfile{NthTropicalMeromorphicFunction} in
\leanfile{Function/PiecewisePolynomial.lean}.  Intrinsic one-sided germs make all
later multiplicities independent of the chosen presentation.
\end{leanproof}

\begin{definition}[Normalized jets and jet matrix]
\paperlabel{def:normalized-jets}
\lean{NthTropicalNevanlinna.normalizedPolynomialJet}
\leanok
For a polynomial $p(X)=\sum_{k=0}^na_kX^k$, its normalized $j$-jet at
$x_0$ is
\[
 J_j(p,x_0)=\frac{p^{(j)}(x_0)}{j!}
 =\sum_{k=j}^n\binom{k}{j}a_kx_0^{k-j}.
\]
Writing $J(p,x_0)=(J_1,\ldots,J_n)^T$ and
$a=(a_1,\ldots,a_n)^T$ gives $J(p,x_0)=\mathcal C_n(x_0)a$, where
$\mathcal C_n(x_0)$ is upper triangular with diagonal entries $1$ and hence
is invertible.  Multiplying its $j$-th row by the appropriate one-sided factor
$\operatorname{sgn}^{j+1}(x_0^\pm)$ gives the signed jet matrix
$\mathcal D_n(x_0^\pm)$ used to recover prescribed multiplicities.
\end{definition}

\begin{definition}[$j$-th roots, poles, and entire functions]
\paperlabel{def:multiplicity}
\lean{NthTropicalNevanlinna.multiplicity,NthTropicalNevanlinna.IsJthRoot,NthTropicalNevanlinna.IsJthPole,NthTropicalNevanlinna.IsTropicalEntire,NthTropicalNevanlinna.IsTropicalNowhereVanishingEntire}
\leanok
\uses{def:nth-meromorphic,def:normalized-jets}
For $1\leq j\leq n$, define
\[
 \omega_f^{(j)}(x)=
 \frac{\operatorname{sgn}^{j+1}(x^+)f^{(j)}(x^+)
 -\operatorname{sgn}^{j+1}(x^-)f^{(j)}(x^-)}{j!},
 \qquad \operatorname{sgn}(0^\pm)=\pm1.
\]
The point $x$ is a $j$-th root if $\omega_f^{(j)}(x)>0$, and a $j$-th
pole if $\omega_f^{(j)}(x)<0$; its multiplicity is
$|\omega_f^{(j)}(x)|$.  The function is \emph{tropical entire} if it has no
poles of any order, and \emph{nowhere-vanishing entire} if it has neither
roots nor poles.
\end{definition}

\begin{example}[2.1 --- Closure failures]
\paperlabel{ex:2.1}
\lean{NthTropicalNevanlinna.example_2_1_counterexamples}
\leanok
\uses{def:multiplicity}
Let $f(x)=\operatorname{sgn}(x)x^2$ and $g(x)=x$.  Both are second-order
tropical entire, but $f\oplusT g$ has a second-order pole of multiplicity
$1$ at $0$, and $f(x+1)$ has a second-order pole of multiplicity $2$ at
$-1$.  Thus translation and max need not preserve higher-order entirety.
\end{example}

\begin{proof}
\leanok
On the two sides of $0$, the maximum $f\oplusT g$ is respectively $-x^2$
and $x$ near the origin.  Substitution into the signed second-jet definition
gives $\omega^{(2)}(0)=-1$.  The translated function has its unique germ
change at $-1$; its two second derivatives have signed jump
$\omega^{(2)}(-1)=-2$.  Both are therefore genuine second-order poles with
the stated multiplicities.
\end{proof}

\begin{leanproof}
\leanfile{example_2_1_counterexamples} in \leanfile{Counterexamples.lean}.
\end{leanproof}

\begin{proposition}[2.2 --- Endpoint mean inequality]
\paperlabel{prop:2.2}
\lean{NthTropicalNevanlinna.entire_midpoint_le_endpointMean,NthTropicalNevanlinna.nowhereVanishingEntire_endpointMean_eq}
\leanok
\uses{def:multiplicity}
If $f$ is $n$-th tropical entire, then for every $r>0$,
\[
                    f(0)\leq\frac{f(r)+f(-r)}2.
\]
If $f$ is nowhere-vanishing entire, equality holds.
\end{proposition}

\begin{proof}
\leanok
Taylor-expand the polynomial pieces from $0$ to $r$ and from $0$ to $-r$.
Adjacent polynomial values cancel, leaving the central polynomial jets and
the weighted sums
$\omega_f^{(j)}(x)(r-|x|)^j$.  Every such term is nonnegative for an entire
function.  If there are no roots either, every jump vanishes, so the two
endpoint expansions add to $2f(0)$.
\end{proof}

\begin{leanproof}
\leanfile{entire_midpoint_le_endpointMean} and
\leanfile{nowhereVanishingEntire_endpointMean_eq} in
\leanfile{Function/EntireMean.lean}.
\end{leanproof}

\begin{example}[Counterexample to the converse of 2.2]
\paperlabel{ex:prop2.2-converse}
\lean{NthTropicalNevanlinna.proposition22_converse_is_false}
\leanok
\uses{prop:2.2}
Define $q(x)=0$ for $|x|\leq1$, $q(x)=-x+1$ for $x>1$, and
$q(x)=-x-1$ for $x<-1$.  Then
$q(r)+q(-r)=2q(0)$ for every $r>0$, but $x=1$ is a genuine first-order
pole.  Thus the endpoint equality does not characterize tropical entirety.
\end{example}

\begin{proof}
\leanok
The two tails are negatives of one another and the middle interval is zero,
which gives the endpoint equality directly.  At $x=1$ the slope changes from
$0$ to $-1$, hence $\omega_q^{(1)}(1)=-1<0$.
\end{proof}

\begin{leanproof}
\leanfile{proposition22_converse_is_false} in \leanfile{Counterexamples.lean}.
\end{leanproof}

\begin{proposition}[2.3 --- Entire quotient decomposition]
\paperlabel{prop:2.3}
\lean{NthTropicalNevanlinna.exists_entire_quotient_decomposition}
\leanok
\uses{def:nth-meromorphic,def:multiplicity}
Every $n$-th tropical meromorphic function $f$ can be written
$f=h\oslashT g=h-g$, where $g$ and $h$ are tropical entire, their orders
$n_1,n_2$ satisfy $\max\{n_1,n_2\}=n$, and they have no common $j$-th root
for $1\leq j\leq n$.
\end{proposition}

\begin{proof}
\leanok
At each genuine singularity prescribe the jump of $g$ to be the positive part
of $-\omega_f$ and put $h=f+g$.  Then
$\omega_g^{(j)}=\max\{-\omega_f^{(j)},0\}$ and
$\omega_h^{(j)}=\max\{\omega_f^{(j)},0\}$.  Polynomial-jet inversion and
continuous spline assembly realize these prescribed jumps globally.  Hence
both functions are entire, have disjoint root data, and $h-g=f$.  Exact
ambient order follows because cancellation of every degree-$n$ coefficient in
both $g$ and $h$ would force the degree-$n$ coefficient of $f$ to vanish.
\end{proof}

\begin{leanproof}
\leanfile{exists_entire_quotient_decomposition} in
\leanfile{Function/Decomposition.lean}; jet inversion and assembly are
encapsulated in \leanfile{PolynomialJet.lean} and \leanfile{SplineAssembly.lean}.
\end{leanproof}

\section{Poisson--Jensen and growth}

Put
\[
 \tau_f^{(j)}(a)=\frac{f^{(j)}(a^+)-f^{(j)}(a^-)}{j!}.
\]

\begin{lemma}[3.1 --- Finite telescoping]
\paperlabel{lem:3.1}
\lean{NthTropicalNevanlinna.endpoint_sub_eq_intrinsic_singularity_sum_of_pos,NthTropicalNevanlinna.sub_leftEndpoint_eq_intrinsic_singularity_sum_of_neg}
\leanok
\uses{def:nth-meromorphic,def:normalized-jets}
Let $f$ be piecewise polynomial of degree at most $n$ on $[-r,r]$, and
$x\in(-r,r)$.  If $x_1<\cdots<x_\nu$ are the genuine derivative-jump points
in $(x,r)$ and $x_{-\mu}<\cdots<x_{-1}$ those in $(-r,x)$, then
\begin{align*}
 f(r)-f(x)
 &=\sum_{j=1}^n\left[
   \frac{f^{(j)}(x^+)}{j!}(r-x)^j
   +\sum_{i=1}^\nu\tau_f^{(j)}(x_i)(r-x_i)^j\right],\\
 f(x)-f(-r)
 &=\sum_{j=1}^n\left[
   (-1)^{j+1}\frac{f^{(j)}(x^-)}{j!}(r+x)^j
   +(-1)^j\sum_{i=1}^\mu\tau_f^{(j)}(x_{-i})(r+x_{-i})^j\right].
\end{align*}
\end{lemma}

\begin{proof}
\leanok
Taylor-expand each polynomial piece at its left endpoint.  Add the increments
along the mesh.  Continuity cancels the zeroth-order terms, and exchanging the
two finite sums collects the jump in each normalized derivative.  The left
formula follows identically after reversing orientation, producing the powers
of $-1$.
\end{proof}

\begin{leanproof}
\leanfile{endpoint_sub_eq_jet_sum}, \leanfile{sub_leftEndpoint_eq_jet_sum}, and their
genuine-singularity wrappers in \leanfile{PoissonJensen/Telescoping.lean} and
\leanfile{PoissonJensen/IntrinsicTelescoping.lean}.
\end{leanproof}

\begin{definition}[Poisson--Jensen kernels]
\paperlabel{def:pj-kernels}
\lean{NthTropicalNevanlinna.symmetricKernel,NthTropicalNevanlinna.antisymmetricKernel,NthTropicalNevanlinna.interactionKernel}
\leanok
For $r>0$, $|x|<r$, and $k\geq0$, define
\[
 B_k(r,x)=\frac{(r+x)^k+(r-x)^k}{2},\qquad
 D_k(r,x)=\frac{(r+x)^k-(r-x)^k}{2},
\]
and for $|y|<r$ define
\[
 E(r,x,y)=r^2-|x-y|r-xy
 =\begin{cases}(r-y)(r+x),&x\leq y,\\
                 (r+y)(r-x),&y\leq x.
   \end{cases}
\]
Thus $B_k\pm D_k=(r\pm x)^k$, $D_k(r,0)=0$, and
$E(r,0,y)=r(r-|y|)$.
\end{definition}

\begin{definition}[Nine-region partition]
\paperlabel{def:pj-regions}
\lean{NthTropicalNevanlinna.RelativeRegion,NthTropicalNevanlinna.InRelativeRegion}
\leanok
For fixed $r>0$ and $|x|<r$, let
\begin{align*}
F_1&=\{y:-r<y<\min(0,x)\},&
F_2&=\{y:\max(0,x)<y<r\},\\
F_3&=\{y:y=x=0\},&
F_4&=\{y:0<y<x\},\\
F_5&=\{y:x<y<0\},&
F_6&=\{y:x<y=0\},\\
F_7&=\{y:y=0<x\},&
F_8&=\{y:0<y=x\},\\
F_9&=\{y:y=x<0\}.&&
\end{align*}
They are disjoint and cover $(-r,r)$; some are empty for a fixed $x$.
\end{definition}

\begin{definition}[Auxiliary jumps]
\paperlabel{def:auxiliary-jumps}
\lean{NthTropicalNevanlinna.auxiliaryOmega,NthTropicalNevanlinna.auxiliaryGamma}
\leanok
\uses{def:multiplicity}
For $1\leq j\leq n$ define
\begin{align*}
\Omega_f^{(j)}(x,y)
 &=\frac{\operatorname{sgn}^{j+1}((y-x)^+)f^{(j)}(y^+)
 -\operatorname{sgn}^{j+1}((y-x)^-)f^{(j)}(y^-)}{j!},\\
\Gamma_f^{(j)}(x,y)
 &=\frac{\operatorname{sgn}^{j}((y-x)^+)f^{(j)}(y^+)
 -\operatorname{sgn}^{j}((y-x)^-)f^{(j)}(y^-)}{j!}.
\end{align*}
In particular $\Omega_f^{(j)}(0,y)=\omega_f^{(j)}(y)$.
\end{definition}

\begin{lemma}[3.2 --- Nine-region reduction]
\paperlabel{lem:3.2}
\lean{NthTropicalNevanlinna.auxiliaryJump_eq_byRegion}
\leanok
\uses{def:pj-regions,def:auxiliary-jumps}
For $y\in F_\ell$, write $L_j(y)=f^{(j)}(y^-)/j!$.  Then
\[
\Omega_f^{(j)}(x,y)=A_{\ell,j}\omega_f^{(j)}(y)
                    +A'_{\ell,j}L_j(y),\qquad
\Gamma_f^{(j)}(x,y)=C_{\ell,j}\omega_f^{(j)}(y)
                    +C'_{\ell,j}L_j(y),
\]
where
\scriptsize
\[
\begin{array}{c|ccccccccc}
\ell&1&2&3&4&5&6&7&8&9\\ \hline
A_{\ell,j}&1&1&1&(-1)^{j+1}&(-1)^{j+1}&1&(-1)^{j+1}&1&(-1)^{j+1}\\
A'_{\ell,j}&0&0&0&0&0&(-1)^{j+1}-1&1-(-1)^{j+1}&1-(-1)^{j+1}&1-(-1)^{j+1}\\
C_{\ell,j}&-1&1&1&(-1)^j&(-1)^{j+1}&1&(-1)^j&1&(-1)^{j+1}\\
C'_{\ell,j}&0&0&2(-1)^{j+1}&0&0&(-1)^{j+1}-1&(-1)^{j+1}-1&1-(-1)^j&1-(-1)^j
\end{array}
\]
\normalsize
\end{lemma}

\begin{proof}
\leanok
Put $u=f^{(j)}(y^+)/j!$, $v=f^{(j)}(y^-)/j!$, and
$\varepsilon=(-1)^{j+1}$.  The multiplicity definition gives
\[
\omega_f^{(j)}(y)=
\begin{cases}u-v,&y>0,\\
\varepsilon(u-v),&y<0,\\
u-\varepsilon v,&y=0.
\end{cases}
\]
The sign of $y-x$ similarly gives
\[
\Omega_f^{(j)}(x,y)=
\begin{cases}u-v,&y>x,\\
\varepsilon(u-v),&y<x,\\
u-\varepsilon v,&y=x,
\end{cases}
\]
while $\Gamma$ is obtained from the same three rows after replacing the
right-side exponent $j+1$ by $j$.  Each region fixes both the sign of $y$ and
the comparison of $y$ with $x$.  Solving the first display for $u-v$ and
substituting into the second gives the four coefficients in the table.  The
cases $F_6,F_7,F_8,F_9$ retain a multiple of $v$ because either $y=0$ or
$y=x$; in all other cases that correction is zero.  Direct substitution for
the nine disjoint regions yields every table entry.
\end{proof}

\begin{leanproof}
\leanfile{auxiliaryJump_eq_byRegion} in
\leanfile{PoissonJensen/AuxiliaryJumps.lean}.  The definitions $B_k,D_k,E$ and the
nine-region partition remain proof-local in \leanfile{Kernels.lean} and
\leanfile{RegionPartition.lean}.
\end{leanproof}

\begin{theorem}[3.3 --- Poisson--Jensen formula]
\paperlabel{thm:3.3}
\lean{NthTropicalNevanlinna.poissonJensen,NthTropicalNevanlinna.jensenFormula}
\leanok
\uses{lem:3.1,def:pj-kernels,lem:3.2}
Let $r>0$, $|x|<r$, and let $S$ be the finite set of genuine singularities of
$f$ in $(-r,r)$.  If $y\in F_\ell$, put
$A_j(x,y)=A_{\ell,j}$ and $C_j(x,y)=C_{\ell,j}$ from Lemma 3.2.  Let
$f_{a^-}$ denote the polynomial germ immediately to the left of $a$, extended
as a polynomial to all real arguments, and let
$s_-(x)=\operatorname{sgn}(x)$ for $x\neq0$ and $s_-(0)=-1$.  Then
\begin{align*}
f(x)={}&\frac{f(r)+f(-r)}2
+\frac{D_n(r,x)}{2B_n(r,x)}\bigl(f(r)-f(-r)\bigr)\\
&-\sum_{j=1}^n\sum_{y\in S}
A_j(x,y)\omega_f^{(j)}(y)E(r,x,y)^j
\frac{B_{n-j}(r,x)}{2B_n(r,x)}\\
&-\sum_{j=1}^n\sum_{y\in S}
C_j(x,y)\omega_f^{(j)}(y)E(r,x,y)^j
\frac{D_{n-j}(r,x)}{2B_n(r,x)}\\
&-\frac{(r+x)^n}{2B_n(r,x)}\bigl(f_{x^-}(r)-f(x)\bigr)
-\frac{(r-x)^n}{2B_n(r,x)}\bigl(f_{x^-}(-r)-f(x)\bigr)\\
&-s_-(x)\frac{(r-|x|)^n}{2B_n(r,x)}
 \bigl(f_{0^-}(r)+f_{0^-}(-r)-2f(0)\bigr).
\end{align*}
At $x=0$ all $D$-terms and germ corrections cancel, giving
\[
 f(0)=\frac{f(r)+f(-r)}2
 -\frac12\sum_{j=1}^n\sum_{\substack{y:\ j\text{-root}\\|y|<r}}
    |\omega_f^{(j)}(y)|(r-|y|)^j
 +\frac12\sum_{j=1}^n\sum_{\substack{z:\ j\text{-pole}\\|z|<r}}
    |\omega_f^{(j)}(z)|(r-|z|)^j.
\]
\end{theorem}

\begin{proof}
\leanok
Multiply the two identities of Lemma 3.1 by $(r+x)^n$ and $(r-x)^n$ and
subtract.  Factor every distance product through $E(r,x,y)$, split symmetric
and antisymmetric powers into $B_k$ and $D_k$, and use Lemma 3.2 to replace
auxiliary jumps by intrinsic multiplicities.  Taylor's formula recombines the
remaining one-sided jets into the two germ corrections.  Division by
$2B_n(r,x)>0$ yields the general formula.  Substitution $x=0$ leaves the
signed root/pole sum above.
\end{proof}

\begin{leanproof}
\leanfile{poissonJensen} and \leanfile{jensenFormula} in
\leanfile{PoissonJensen/Main.lean}.  The public theorem constructs the needed
finite interval presentation from the global function.
\end{leanproof}

\begin{example}[3.4 --- Multiplicity table]
\paperlabel{ex:3.4}
\lean{NthTropicalNevanlinna.example34_complete_multiplicity_table,NthTropicalNevanlinna.example34_jensen_numerical_identity}
\leanok
\uses{thm:3.3}
On $(-3,3)$ let
\[
f(x)=\begin{cases}
-2x^3-17x^2-46x-40,&-3<x\leq-2,\\
2x^2+6x+4,&-2<x\leq-1,\\
x^2+2x+1,&-1<x\leq1,\\
-x+5,&1<x\leq2,\\
3x^3-15x^2+18x+3,&2<x<3.
\end{cases}
\]
Its nonzero signed multiplicities are
\[
\begin{array}{c|rrrrr}
x&-2&-1&0&1&2\\ \hline
\omega_f^{(1)}&&-2&&-5&-5\\
\omega_f^{(2)}&-7&1&2&-1&3\\
\omega_f^{(3)}&2&&&&3
\end{array}
\]
and for $2<r<3$ the third-order Jensen formula evaluates to $f(0)=1$.
\end{example}

\begin{proof}
\leanok
Differentiate the two polynomial germs at each of
$-2,-1,0,1,2$, divide the $j$-th derivative jump by $j!$, and insert the
one-sided sign factors.  This gives exactly the displayed table; blank
entries are zero.  Separating its positive and negative entries into roots
and poles and substituting their weights $(r-|x|)^j/2$ into the Jensen formula
cancels every cubic, quadratic, and linear term, leaving $1=f(0)$.
\end{proof}

\begin{leanproof}
\leanfile{example34_complete_multiplicity_table} in
\leanfile{Counterexamples.lean}.
\end{leanproof}

\begin{definition}[Nevanlinna quantities]
\paperlabel{def:nevanlinna-quantities}
\lean{NthTropicalNevanlinna.proximity,NthTropicalNevanlinna.maxPlusCounting,NthTropicalNevanlinna.integratedCounting,NthTropicalNevanlinna.partialIntegratedCounting,NthTropicalNevanlinna.characteristic,NthTropicalNevanlinna.growthOrder,NthTropicalNevanlinna.hyperOrder,NthTropicalNevanlinna.integratedCounting_eq_paperIteratedIntegral}
\leanok
\uses{def:multiplicity}
For $r>0$ define
\begin{align*}
 m(r,f)&=\frac{\max\{f(r),0\}+\max\{f(-r),0\}}2,\\
 N^{(j)}(r,f)&=\frac12\sum_{z:\ j\text{-pole},\ |z|<r}
 |\omega_f^{(j)}(z)|(r-|z|)^j,\\
 T(r,f)&=m(r,f)+\sum_{j=1}^nN^{(j)}(r,f).
\end{align*}
For a region $A\subseteq\R$, $N_A^{(j)}(r,f)$ is the same finite sum
restricted to $z\in A$.  If $n^{(j)}(t,f)$ counts $j$-th poles in $(-t,t)$,
then
\[
 N^{(j)}(r,f)=\frac12\int_{[0,r]^j}
 n^{(j)}\!\left(\min_s t_s,f\right)\,dt_1\cdots dt_j.
\]
Finally
\[
 \rho(f)=\limsup_{r\to\infty}\frac{\log T(r,f)}{\log r},\qquad
 \rho_2(f)=\limsup_{r\to\infty}\frac{\log\log T(r,f)}{\log r}.
\]
\end{definition}

\begin{proof}[Sum--integral identity]
\leanok
Expand the counting function into the finite sum of indicators of boxes
$[|z|,r]^j$.  Tonelli reduces the integral to a finite sum, and the box
attached to $z$ has volume $(r-|z|)^j$.
\end{proof}

\begin{leanproof}
Definitions \leanfile{proximity}, \leanfile{maxPlusCounting},
\leanfile{integratedCounting}, \leanfile{partialIntegratedCounting}, and
\leanfile{characteristic}; the bridge and integral identity are
\leanfile{countingBoxIntegrand_eq_maxPlusCounting_minimum},
\leanfile{integratedCounting_eq_iteratedIntegral}, and
\leanfile{integratedCounting_eq_paperIteratedIntegral} in
\leanfile{Nevanlinna/Counting.lean}.  Growth invariants are in
\leanfile{Nevanlinna/Growth.lean}.
\end{leanproof}

\begin{theorem}[3.5 --- Jensen identity]
\paperlabel{thm:3.5}
\lean{NthTropicalNevanlinna.characteristic_eq_neg_add_at_zero}
\leanok
\uses{thm:3.3,def:nevanlinna-quantities}
For every $r>0$,
\[
                         T(r,f)=T(r,-f)+f(0).
\]
\end{theorem}

\begin{proof}
\leanok
Split the signed multiplicity sum in the Jensen formula into positive roots
and negative poles.  Replacing $f$ by $-f$ exchanges the two sets, while
$\max(a,0)-\max(-a,0)=a$ converts the endpoint term into
$(f(r)+f(-r))/2$.  Rearrangement gives the identity.
\end{proof}

\begin{leanproof}
\leanfile{characteristic_eq_neg_add_at_zero} in
\leanfile{Nevanlinna/Jensen.lean}.
\end{leanproof}

\begin{definition}[3.6--3.7 --- Tropical hyper-exponentials]
\paperlabel{def:hyper-exponential}
\lean{NthTropicalNevanlinna.tropicalHyperExponential}
\leanok
For $\alpha>1$, the first-order tropical hyper-exponential is
\[
e_\alpha(x)=\alpha^{\lfloor x\rfloor}
\left(x-\lfloor x\rfloor+\frac1{\alpha-1}\right).
\]
For $n\geq1$ define
\begin{align*}
e_{n,\alpha}(x)={}&
\operatorname{sgn}^{n+1}(x^+)\alpha^{\lfloor x\rfloor}
\bigl(x^n-\lfloor x\rfloor^n\bigr)\\
&+\sum_{i=-\infty}^{\lfloor x\rfloor-1}
\operatorname{sgn}^{n+1}(i^+)\alpha^i
\bigl((i+1)^n-i^n\bigr).
\end{align*}
The negative infinite tail is interpreted as its absolutely convergent
geometric-polynomial series.  For $n=1$ this reduces to $e_\alpha$.
\end{definition}

\begin{proposition}[3.8 --- Hyper-exponential growth]
\paperlabel{prop:3.8}
\lean{NthTropicalNevanlinna.proposition_3_8}
\leanok
\uses{def:hyper-exponential,def:nevanlinna-quantities}
For $\alpha>1$, the generalized tropical hyper-exponential $e_{n,\alpha}$ is
a nonnegative, nondecreasing, $n$-th tropical entire function and
$\rho_2(e_{n,\alpha})=1$.
\end{proposition}

\begin{proof}
\leanok
Reindex the negative tail by $k=-i$; it is bounded by a finite linear
combination of series $\sum_{k\geq0}k^q\alpha^{-k}$ and therefore converges
absolutely.  At an integer $m$, the left polynomial contributes precisely
the $i=m-1$ term missing from the right tail, so the two one-sided limits
agree.  On a nonintegral interval the derivative is
$n\alpha^{\lfloor x\rfloor}|x|^{n-1}\geq0$, proving monotonicity; the limit at
$-\infty$ and monotonicity give nonnegativity.

At an integer, direct normalized-jet computation gives nonnegative
$\omega_{e_{n,\alpha}}^{(j)}$ for every $1\leq j\leq n$, while away from the
integers every multiplicity is zero.  Hence the function is tropical entire.
For large positive $r$, its endpoint value is bounded above by
$C(1+r)^n\alpha^r$ and below on a fixed terminal part of each unit interval
by $c\alpha^{r-1}$.  Jensen and nonnegativity compare these bounds with the
characteristic.  Taking two logarithms and dividing by $\log r$ squeezes the
limsup between $1$ and $1$, so $\rho_2(e_{n,\alpha})=1$.
\end{proof}

\begin{leanproof}
\leanfile{tropicalHyperExponential} and \leanfile{proposition_3_8} in
\leanfile{Nevanlinna/Examples.lean}.
\end{leanproof}

\section{Logarithmic differences}

\begin{example}[4.1 --- Failure of monotonicity and convexity]
\paperlabel{ex:4.1}
\lean{NthTropicalNevanlinna.example41_characteristic_not_monotone,NthTropicalNevanlinna.example41_characteristic_not_convex}
\leanok
\uses{def:nevanlinna-quantities}
For every $m\in\N\cup\{0\}$, define
\[
 f(x)=\begin{cases}
 -(x-m)(x-m-1),&m\leq x\leq m+1,\\
 (x+m)(x+m+1),&-m-1\leq x\leq-m.
 \end{cases}
\]
Then
\begin{align*}
 m(r,f)&=-\frac12(r-\lfloor r\rfloor)(r-\lfloor r\rfloor-1),\\
 \sum_{j=1}^2N^{(j)}(r,f)&=N^{(1)}(r,f)
   =\frac12\sum_{i=1}^{\lfloor r\rfloor}2(r-i),\\
 T(r,f)&=-\frac12r^2+\left(2\lfloor r\rfloor+\frac12\right)r
          -\lfloor r\rfloor(\lfloor r\rfloor+1).
\end{align*}
In particular, $T(r,f)$ is neither nondecreasing nor convex.
\end{example}

\begin{proof}
\leanok
On each unit interval the endpoint values give the displayed proximity term.
The first derivatives jump down by $2$ at every nonzero integer of the
appropriate sign, while the normalized second jets agree; hence only the
first-order counting term remains and summing its weights gives the second
formula.  Their sum is the displayed expression for $T$.  On $0\leq r\leq1$
it reduces to $(-r^2+r)/2$: it decreases after $r=1/2$ and has second
derivative $-1$.  Thus it is neither nondecreasing nor convex.
\end{proof}

\begin{leanproof}
\leanfile{example41_characteristic_not_monotone} and
\leanfile{example41_characteristic_not_convex} in \leanfile{Counterexamples.lean}.
\end{leanproof}

\begin{definition}[4.2 --- Well-defined polynomials]
\paperlabel{def:4.2}
\lean{NthTropicalNevanlinna.IsWellDefinedPolynomial,NthTropicalNevanlinna.IsWellDefinedNthTropicalMeromorphicFunction}
\leanok
\uses{def:nth-meromorphic}
A polynomial is \emph{well-defined} if, after subtracting its value at zero,
it contains only odd powers when its degree is odd or only even powers when
its degree is even, and all nonzero coefficients have one common weak sign.
An $n$-th tropical meromorphic function is well-defined when each polynomial
piece is well-defined.
\end{definition}

\begin{leanproof}
\leanfile{IsWellDefinedPolynomial} and
\leanfile{IsWellDefinedNthTropicalMeromorphicFunction} in
\leanfile{LogDerivative/WellDefined.lean}.
\end{leanproof}

\begin{lemma}[4.3 --- Radial derivative signs]
\paperlabel{lem:4.3}
\lean{NthTropicalNevanlinna.wellDefinedPolynomial_iteratedDeriv_commonSign}
\leanok
\uses{def:4.2}
Let $p$ be a well-defined polynomial of degree $n$, $\delta\in\{-1,1\}$,
and $r>0$.  For all $1\leq\ell\leq n$, the derivatives
$d^\ell p(\delta r)/dr^\ell$ have one common weak sign.
\end{lemma}

\begin{proof}
\leanok
If $n=2m$, then
\[
 \frac{d^\ell}{dr^\ell}p(\delta r)=
 \sum_{2j\geq\ell}a_{2j}\frac{(2j)!}{(2j-\ell)!}r^{2j-\ell}.
\]
Every weight is nonnegative, so the sum has the common coefficient sign.  If
$n=2m-1$, the same calculation gives
\[
 \frac{d^\ell}{dr^\ell}p(\delta r)=
 \delta\sum_{2j-1\geq\ell}a_{2j-1}
 \frac{(2j-1)!}{(2j-1-\ell)!}r^{2j-1-\ell},
\]
whose sign is the coefficient sign multiplied by the fixed $\delta$.
\end{proof}

\begin{leanproof}
\leanfile{wellDefinedPolynomial_iteratedDeriv_commonSign} in
\leanfile{LogDerivative/WellDefined.lean}.
\end{leanproof}

\begin{definition}[Well-defined $n$-th tropical meromorphic function]
\paperlabel{def:well-defined-function}
\lean{NthTropicalNevanlinna.IsWellDefinedNthTropicalMeromorphicFunction}
\leanok
\uses{def:4.2}
Such a function is well-defined when every polynomial germ in one, hence any,
intrinsic piecewise-polynomial realization is well-defined in Definition 4.2.
\end{definition}

\begin{lemma}[4.4 --- Characteristic monotonicity]
\paperlabel{lem:4.4}
\lean{NthTropicalNevanlinna.characteristicLemma44}
\leanok
\uses{lem:4.3,def:nevanlinna-quantities,def:well-defined-function}
For a well-defined $n$-th tropical meromorphic function $f$, every left radial
derivative $T^{(\ell)}(r^-,f)$, $0\leq\ell\leq n$, is nonnegative and
nondecreasing.  In particular $T(\cdot,f)$ is nonnegative, nondecreasing, and
convex on $(0,\infty)$.
\end{lemma}

\begin{proof}
\leanok
Put $h_\delta(s)=f(\delta s)$.  At an event radius $a>0$, write
$u_\delta^\pm=h_\delta^{(\ell)}(a^\pm)$ and let $\chi_\delta^\pm$ indicate
that $h_\delta$ is positive on the corresponding side.  The multiplicity
definition gives
\[
 u_\delta^+-u_\delta^-=\ell!\,\omega_f^{(\ell)}(\delta a).
\]
The combined jump of the proximity and pole-counting terms in the
$\ell$-th derivative is
\[
 J_{\delta,\ell}(a)=\frac12\left(
 \chi_\delta^+u_\delta^+-\chi_\delta^-u_\delta^-
 +\max\{u_\delta^--u_\delta^+,0\}\right).
\]
If $h_\delta(a)>0$, this is
$\max\{u_\delta^+-u_\delta^-,0\}/2$; if $h_\delta(a)<0$, it is
$\max\{u_\delta^--u_\delta^+,0\}/2$.  If $h_\delta(a)=0$, Lemma 4.3 gives
$\chi_\delta^+u_\delta^+\geq0$ and
$-\chi_\delta^-u_\delta^-\geq0$.  Thus every derivative jump is
nonnegative.

Near the origin, Taylor expansion gives
\[
 T(r,f)=T(0^+,f)+\sum_{j=1}^nc_jr^j,\qquad
 c_j=\frac{\chi_+a_j+\chi_-b_j+max\{-(a_j+b_j),0\}}{2j!},
\]
where $a_j=d^jf(r)/dr^j|_{0^+}$ and
$b_j=d^jf(-r)/dr^j|_{0^+}$.  If $f(0)$ is positive or negative this reduces
respectively to the positive or negative part of $a_j+b_j$; if $f(0)=0$,
Lemma 4.3 controls the two indicated terms separately.  Hence $c_j\geq0$.

Between locally finite event radii, $T$ is a polynomial of degree at most
$n$.  Its $n$-th derivative is therefore constant there and has only
nonnegative jumps, so it is nonnegative and nondecreasing.  Descending in
$\ell$, integrate between events and use the nonnegative jumps and the
initial coefficients.  This proves the assertion for every $\ell$; the cases
$\ell=0,1$ give monotonicity and convexity.
\end{proof}

\begin{leanproof}
\leanfile{characteristicLemma44} in
\leanfile{LogDerivative/CharacteristicMonotonicity.lean}.
\end{leanproof}

\begin{lemma}[4.5 --- Pointwise shift estimate]
\paperlabel{lem:4.5}
\lean{NthTropicalNevanlinna.pointwiseShiftEstimate_of_wellDefined}
\leanok
\uses{lem:3.1,thm:3.5,lem:4.4}
If $f$ is well-defined, $\alpha>1$, $c\neq0$, $\delta\in\{-1,1\}$, and
\[
 r>\max\left\{2|c|,\frac{3-\alpha}{\alpha-1}|c|\right\},
\]
then
\[
 |f(\delta r+c)-f(\delta r)|\leq
 \frac{32|c|}{(\alpha-1)(r+|c|)}
 \left(T(\alpha(r+|c|),f)+\frac{|f(0)|}{2}\right).
\]
\end{lemma}

\begin{proof}
\leanok
Set $\rho=(\alpha+1)(r+|c|)/2$.  Apply Lemma 3.1 on the interval from
$\delta r$ to $\delta r+c$ and again out to $\delta\rho$.  Common-sign jets
permit the short increments to be bounded by the long increments.  Split
each singularity contribution into root and pole parts, dominate the partial
counts by the full counts at $\rho$, and use the Jensen identity to replace
the characteristic of $-f$.  Monotonicity from Lemma 4.4 and the definition
of $\rho$ give the displayed constant.
\end{proof}

\begin{leanproof}
\leanfile{pointwiseShiftEstimate_of_wellDefined} in
\leanfile{LogDerivative/ShiftEstimate.lean}.
\end{leanproof}

\begin{lemma}[4.6 --- Growth shift]
\paperlabel{lem:4.6}
\lean{NthTropicalNevanlinna.growthShiftLemma}
\leanok
\uses{def:nevanlinna-quantities}
Let $T:[0,\infty)\to[0,\infty)$ be continuous and nondecreasing with
hyper-order $\rho_2<1$.  If $c>0$ and $0<\sigma<1-\rho_2$, then outside a set
of finite logarithmic measure,
\[
                       T(r+c)=T(r)+o\!\left(\frac{T(r)}{r^\sigma}\right).
\]
\end{lemma}

\begin{proof}
\leanok
Choose $\widetilde\sigma$ with
$\sigma<\widetilde\sigma<1-\rho_2$.  If $T$ eventually vanishes there is
nothing to prove; otherwise take $R>\max\{1,c\}$ with $T(R)>0$.  For
$\eta>0$ set
\[
 F_\eta=\left\{r\geq R:
 \frac{T(r+c)-T(r)}{T(r)}r^{\widetilde\sigma}\geq\eta\right\}.
\]
We prove that $F_\eta$ has finite logarithmic measure.  Suppose instead that
$\int_{F_\eta}dt/t=\infty$.  Continuity makes $F_\eta$ closed and it must be
unbounded.  Recursively choose
\[
 r_0=\min F_\eta,\qquad
 r_{m+1}=\min(F_\eta\cap[r_m+c,\infty)).
\]
Then
\[
 F_\eta\subseteq\bigcup_{m\geq0}[r_m,r_m+c],\qquad
 T(r_{m+1})\geq
 \left(1+\frac{\eta}{r_m^{\widetilde\sigma}}\right)T(r_m).       \tag{4.6.1}
\]

For every $\varepsilon>0$, infinitely many $m$ satisfy
$r_m<m^{1+\varepsilon}$.  Otherwise, from the covering and
$\log(1+x)\leq x$,
\[
 \int_{F_\eta}\frac{dt}{t}
 \leq C+\sum_{m\gg1}\log\frac{r_m+c}{r_m}
 \leq C+c\sum_{m\gg1}m^{-1-\varepsilon}<\infty,
\]
a contradiction.  Choose $\varepsilon$ so that
$1/(1+\varepsilon)>\widetilde\sigma$ and pass to such a subsequence.
Iterating (4.6.1), using $r_k\leq r_m$, and then
$\log(1+x)\geq x/2$ for small $x>0$, gives
\[
 \log T(r_m)\geq C_0+
 \frac{\eta m}{2r_m^{\widetilde\sigma}}.
\]
The right side tends to infinity, and a second logarithm yields
\[
 \frac{\log\log T(r_m)}{\log r_m}
 \geq\frac{\log m}{\log r_m}-\widetilde\sigma+o(1)
 \geq\frac1{1+\varepsilon}-\widetilde\sigma+o(1).
\]
Thus $\rho_2\geq1/(1+\varepsilon)-\widetilde\sigma$; letting
$\varepsilon\downarrow0$ contradicts
$\widetilde\sigma<1-\rho_2$.  Hence $F_\eta$ has finite logarithmic measure.

Take $E=F_1$.  For $r\notin E$,
\[
 0\leq T(r+c)-T(r)<\frac{T(r)}{r^{\widetilde\sigma}}
 =\frac{T(r)}{r^\sigma}r^{-(\widetilde\sigma-\sigma)},
\]
and the final factor tends to zero.  This is the claimed estimate.
\end{proof}

\begin{leanproof}
\leanfile{growthShiftLemma} in \leanfile{LogDerivative/GrowthShiftProof.lean}.
\end{leanproof}

\begin{theorem}[4.7 --- Pointwise logarithmic derivative]
\paperlabel{thm:4.7}
\lean{NthTropicalNevanlinna.pointwiseLogarithmicDerivative}
\leanok
\uses{lem:4.5,lem:4.6}
If $f$ is well-defined, $c\neq0$, $\rho_2(f)<1$, and
$0<\sigma<1-\rho_2(f)$, then there is a set $E$ of finite logarithmic
measure such that, for both $\delta=\pm1$,
\[
 |f(\delta r+c)-f(\delta r)|
 =o\!\left(\frac{T(r,f)}{r^\sigma}\right)
 \quad(r\to\infty, r\notin E).
\]
\end{theorem}

\begin{proof}
\leanok
We first prove the power-step estimate used below.  If $T$ is nondecreasing,
has hyper-order $\rho_2$, and $0<q<1-\rho_2$, then outside a set $E_q$ of
finite logarithmic measure,
\[
                         T(r+r^q)\leq2T(r).                       \tag{4.7.1}
\]
Choose $\rho_2<\beta<1-q$.  Eventually $|\log T(x)|\leq x^\beta$ after
discarding the identically-zero case.  Put $U=\log T$ and let
$F=\{x:U(x+x^q)-U(x)\geq\log2\}$.  On the exponential annulus
$[a_m,b_m]=[Re^m,Re^{m+1}]$, set $H_m=b_m^q$.  Monotonicity and the
sliding-slope identity give
\begin{align*}
 \int_{F\cap[a_m,b_m]}\frac{dx}{x}
 &\leq\frac{H_m}{(\log2)a_m}
       \bigl(U(b_m+H_m)-U(a_m)\bigr)\\
 &\leq C a_m^{q+\beta-1}.
\end{align*}
The geometric series converges because $q+\beta-1<0$.  Hence $F$ has finite
logarithmic measure, and exponentiation outside $F$ proves (4.7.1).

Apply this to $T(r,f)$ with
$q=(\sigma+1-\rho_2)/2$, so $\sigma<q<1-\rho_2$.  For large
$r\notin E_q$, let $s=r+|c|$ and $A_r=(r+r^q)/s$.  Then
\[
 (A_r-1)s=r^q-|c|,\qquad A_rs=r+r^q.
\]
Eventually $A_r>1$ and the radius hypothesis of Lemma 4.5 holds.  Lemma 4.5
and (4.7.1) therefore imply, simultaneously for $\delta=\pm1$,
\[
 |f(\delta r+c)-f(\delta r)|
 \leq\frac{32|c|}{r^q-|c|}
 \left(T(r+r^q,f)+\frac{|f(0)|}{2}\right)
 \leq C\frac{T(r,f)}{r^q}.
\]
Here, if $T(R_0,f)>0$, monotonicity absorbs $|f(0)|$ into $T(r,f)$.
Since $q>\sigma$, the last quantity is
$r^{\sigma-q}T(r,f)/r^\sigma=o(T(r,f)/r^\sigma)$.

If $T(r,f)\equiv0$, apply Lemma 4.5 with any sufficiently large constant
$A$ in place of $\alpha$ and let $A\to\infty$; this forces every displayed
shift difference to be zero.  Thus the same conclusion holds in the
degenerate case, and one common exceptional set works for both signs.
\end{proof}

\begin{corollary}[4.8 --- Logarithmic derivative lemma]
\paperlabel{cor:4.8}
\lean{NthTropicalNevanlinna.logarithmicDerivativeProximity}
\leanok
\uses{thm:4.7}
Under the same hypotheses,
\[
 m\bigl(r,f(x+c)\oslashT f(x)\bigr)
 =o\!\left(\frac{T(r,f)}{r^\sigma}\right)
\]
outside a set of finite logarithmic measure.
\end{corollary}

\begin{proof}
\leanok
At the two radial endpoints, $u^+\leq|u|$.  Average the two estimates from
Theorem 4.7.
\end{proof}

\begin{leanproof}
\leanfile{pointwiseLogarithmicDerivative} and
\leanfile{logarithmicDerivativeProximity} in \leanfile{LogDerivative/Main.lean}.
\end{leanproof}

\section{Holomorphic curves}

\begin{definition}[Tropical holomorphic curves]
\paperlabel{def:curve-objects}
\lean{NthTropicalNevanlinna.TropicalHolomorphicCurveRepresentation,NthTropicalNevanlinna.NthTropicalHolomorphicCurve,NthTropicalNevanlinna.cartanCharacteristic}
\leanok
\uses{def:nth-meromorphic,def:multiplicity}
An $n$-th tropical holomorphic curve is a projective class
$f=[f_0:\cdots:f_m]$ of tropical entire coordinates whose maximum order is
$n$.  A representation is reduced when the coordinates have no common
$j$-th root.  With $F(x)=\max_i f_i(x)$, define
\[
                  T_f(r)=\frac{F(r)+F(-r)}2-F(0).
\]
\end{definition}

\begin{leanproof}
\leanfile{TropicalHolomorphicCurveRepresentation},
\leanfile{NthTropicalHolomorphicCurve}, and \leanfile{cartanCharacteristic} in
\leanfile{Curves/Basic.lean} and \leanfile{Curves/Characteristic.lean}.
\end{leanproof}

\begin{proposition}[5.2 --- Representation invariance]
\paperlabel{prop:5.2}
\lean{NthTropicalNevanlinna.cartanCharacteristic_eq_of_reduced_projectiveValue_eq_arbitraryOrders}
\leanok
\uses{prop:2.2,def:curve-objects}
The Cartan characteristic is independent of the reduced representation.
\end{proposition}

\begin{proof}
\leanok
Two reduced representatives differ by a common function $\lambda$.  A root
or pole of $\lambda$ would create a common root after adding or subtracting
it from every coordinate, contradicting reducedness.  Thus $\lambda$ is
nowhere-vanishing entire.  Proposition 2.2 gives
$\bigl(\lambda(r)+\lambda(-r)\bigr)/2=\lambda(0)$, so the common additive
change cancels from the defining expression for $T_f$.
\end{proof}

\begin{proposition}[5.3 --- Projective-pair characteristic]
\paperlabel{prop:5.3}
\lean{NthTropicalNevanlinna.cartanCharacteristic_projectivePair_eq_characteristic_sub_arbitraryOrders}
\leanok
\uses{prop:2.3,prop:5.2,thm:3.5,def:curve-objects}
Let $g$ have order $n\geq1$, with $g=g_1-g_0$, where $g_0,g_1$ are entire
and have no common roots of any order.  Their maximum order need not equal
$n$.  For $r>0$,
\[
 T_{[g_0:g_1]}(r)=T(r,g)-g^+(0).
\]
\end{proposition}

\begin{proof}
\leanok
Proposition 2.3 gives another reduced decomposition $g=h_1-h_0$ with
$\max(\deg h_0,\deg h_1)=n$.  The two coordinate pairs differ by the common
function $g_0-h_0$, so $[g_0:g_1]=[h_0:h_1]$.  Proposition 5.2 gives equal
Cartan characteristics, even if the two representations have different
maximum orders: its multiplicity argument applies through the larger order,
and all higher multiplicities vanish.

For the chosen pair, $\max(h_0,h_1)=h_0+g^+$.  Jensen and the absence of
common roots identify the root counts of $h_0$ with the pole counts of $g$,
giving
\[
 T_{[h_0:h_1]}(r)
 =m(r,g)+\sum_{j=1}^{n}N^{(j)}(r,g)-g^+(0)
 =T(r,g)-g^+(0).
\]
Transporting this equality back to the original pair proves the claim.
\end{proof}

\begin{proposition}[5.4 --- Coordinate quotient bound]
\paperlabel{prop:5.4}
\lean{NthTropicalNevanlinna.characteristic_coordinateQuotient_le_upToConstant}
\leanok
\uses{prop:2.2,prop:5.3}
For every coordinate quotient of a reduced curve,
\[
                       T(r,f_i-f_\ell)\leq T_f(r)+O(1).
\]
\end{proposition}

\begin{proof}
\leanok
Construct the common-root part $u$ of $f_i$ and $f_\ell$ by taking the
minimum of each positive multiplicity and fixing $u(0)=f_\ell(0)$.  Then
$h_i=f_i-u$ and $h_\ell=f_\ell-u$ form a reduced pair.  Proposition 5.3 and
$\max(h_i,h_\ell)=\max(f_i,f_\ell)-u$ give
\[
 T(r,f_i-f_\ell)
 \leq\frac12\sum_{\delta=\pm1}
   \bigl(\max_k f_k(\delta r)-u(\delta r)\bigr)+O(1).
\]
Because $u$ is entire, Proposition 2.2 bounds its negative endpoint mean by
a constant.  The right side is therefore at most $T_f(r)+O(1)$.
\end{proof}

\begin{leanproof}
\leanfile{cartanCharacteristic_eq_of_reduced_projectiveValue_eq_arbitraryOrders},
\leanfile{cartanCharacteristic_projectivePair_eq_characteristic_sub_arbitraryOrders}, and
\leanfile{characteristic_coordinateQuotient_le_upToConstant} in
\leanfile{Curves/Characteristic.lean}.
\end{leanproof}

\begin{definition}[Homogeneous and Fermat tropical polynomials]
\paperlabel{def:homogeneous-polynomials}
\lean{NthTropicalNevanlinna.HomogeneousTropicalPolynomial,NthTropicalNevanlinna.HomogeneousTropicalPolynomial.eval,NthTropicalNevanlinna.HomogeneousTropicalPolynomial.curveComposition,NthTropicalNevanlinna.TropicalHomogeneousFermatPolynomial}
\leanok
\uses{def:curve-objects}
A homogeneous tropical polynomial of degree $d$ is a finite maximum of
monomials
\[
 P(x_0,\ldots,x_m)=\max_{|I|=d}
   \left(\alpha_I+\sum_{k=0}^m I_kx_k\right),
 \qquad |I|=\sum_kI_k=d.
\]
Composition is
$(P\circ f)(x)=\max_{|I|=d}(\alpha_I+\sum_kI_kf_k(x))$.  A tropical Fermat
polynomial retains only the pure terms $\max_i(\alpha_i+dx_i)$.
\end{definition}

\begin{example}[5.5 --- Counting dominates curve growth]
\paperlabel{ex:5.5}
\lean{NthTropicalNevanlinna.example55_characteristic_isLittleO_integratedCounting}
\leanok
\uses{def:homogeneous-polynomials,def:nevanlinna-quantities}
Let $P(x_0,x_1)=x_0\oplusT x_1$, $f_0(x)=\operatorname{sgn}(x)x^2$, and
\[
 f_1(x)=\begin{cases}
 -\frac12,&x\leq0,\\
 (4k+2)x-2k(2k+2)-\frac12,&2k\leq x<2k+2,\quad k\geq0.
 \end{cases}
\]
Then $T_f(r)=\frac12r^2(1+o(1))$, whereas
\[
 \sum_{j=1}^2N^{(j)}(r,P\circ f)\geq
 \frac12\sum_{i=1}^{\lfloor r/2\rfloor}(r-2i)^2
 =\frac{r^3}{12}(1+o(1)).
\]
Consequently $T_f=o(\sum_jN^{(j)}(r,P\circ f))$.
\end{example}

\begin{proof}
\leanok
On $[2k,2k+2)$, solving $x^2=f_1(x)$ gives
$x=2k+1\pm\sqrt2/2$.  At the lower crossing the maximum switches from the
affine branch to the quadratic branch and its normalized second-jet jump is
$-1$, so this is a genuine second-order pole of multiplicity $1$.  Summing
their integrated weights gives the displayed lower bound.  The maximum of
the two coordinates has quadratic endpoint growth $r^2+o(r^2)$ and vanishes
at the central normalization up to a constant, giving
$T_f(r)=r^2/2+o(r^2)$.  Cubic growth dominates quadratic growth.
\end{proof}

\begin{leanproof}
\leanfile{example55_characteristic_isLittleO_integratedCounting} in
\leanfile{Counterexamples.lean}.
\end{leanproof}

\begin{theorem}[5.6 --- Tropical polynomial second main theorem]
\paperlabel{thm:5.6}
\lean{NthTropicalNevanlinna.secondMain_homogeneousTropicalPolynomial}
\leanok
\uses{thm:3.5,def:homogeneous-polynomials,def:curve-objects}
Let $P$ be homogeneous of degree $d>0$ and contain every pure power with a
finite coefficient.  For a nonconstant reduced curve,
\[
 T_f(r)=\frac1d\left(
 \sum_{j=1}^nN^{(j)}(r,-P\circ f)
 -\sum_{j=1}^nN^{(j)}(r,P\circ f)\right)+O(1).
\]
\end{theorem}

\begin{proof}
\leanok
The pure terms and the maximum coefficient give uniform bounds
\[
 \gamma+d\max_i f_i(x)\leq P\circ f(x)
 \leq\beta+d\max_i f_i(x).
\]
Apply Jensen to $P\circ f$:
\[
 \sum_jN^{(j)}(r,-P\circ f)-\sum_jN^{(j)}(r,P\circ f)
 =\frac{P\circ f(r)+P\circ f(-r)}2-P\circ f(0).
\]
Average the two pointwise bounds.  The right side equals $dT_f(r)+O(1)$;
division by $d$ proves the formula.
\end{proof}

\begin{corollary}[5.7 --- First-order Fermat case]
\paperlabel{cor:5.7}
\lean{NthTropicalNevanlinna.secondMain_firstOrder_tropicalFermat}
\leanok
\uses{thm:5.6}
For a nonconstant reduced first-order curve and a tropical Fermat polynomial
of degree $d>0$,
\[
 T_f(r)=\frac1dN^{(1)}(r,-P\circ f)+O(1).
\]
\end{corollary}

\begin{proof}
\leanok
Every first-order entire coordinate is convex piecewise linear; multiplication
by $d>0$, addition of constants, and a finite maximum preserve convexity.
Thus $P\circ f$ is tropical entire and its pole count vanishes.  Apply
Theorem 5.6.
\end{proof}

\begin{leanproof}
\leanfile{secondMain_homogeneousTropicalPolynomial} and
\leanfile{secondMain_firstOrder_tropicalFermat} in
\leanfile{Curves/SecondMain.lean}.  Composition closure is constructed, not
assumed.
\end{leanproof}

\begin{theorem}[5.8 --- Ordinary Fermat estimate]
\paperlabel{thm:5.8}
\lean{NthTropicalNevanlinna.secondMain_ordinaryFermat}
\leanok
\uses{def:curve-objects,def:nevanlinna-quantities}
Let $f$ be a nonconstant reduced first-order curve and
$\mathcal P(x)=\sum_{i=0}^m\alpha_ix_i^n$ with $n\geq1$ and $\alpha_i>0$.
If $r/T_f(r)\to0$, then
\[
 \theta T_f(r)^n
 \leq \sum_{j=1}^nN^{(j)}(r,-\mathcal P\circ f)+o(T_f(r)^n)
 \leq \Theta T_f(r)^n,
\]
where $\theta=\min_i\alpha_i$ and
$\Theta=2^{n-1}\sum_i\alpha_i$.
\end{theorem}

\begin{proof}
\leanok
On a linear germ $g(x)=kx+b$,
\[
 \frac{(g^n)^{(j)}}{j!}=\binom njg^{,n-j}k^j,
\]
so every singularity of $f_i^n$ lies at a singularity of $f_i$ or at the
origin.  A first-order entire coordinate is convex and has one of two tail
types: its two slopes eventually point outward, or all slopes have one weak
sign.  In the first case all sufficiently remote jumps of $f_i^n$ are roots,
and the compact remainder is $O(r^n)$.  In the second, at most one sign
crossing occurs; comparison with the affine tail and first-order Jensen gives
a total remote contribution $O(r^n)$.  Since $r=o(T_f)$,
\[
                  \sum_{j=1}^nN^{(j)}(r,f_i^n)=o(T_f(r)^n).
\]
Subadditivity of pole counting and positive homogeneity give the same estimate
for $\mathcal P\circ f$.

Jensen now yields
\[
 \sum_{j=1}^nN^{(j)}(r,-\mathcal P\circ f)
 =\frac12\sum_{\delta=\pm1}\mathcal P\circ f(\delta r)
   +o(T_f(r)^n).
\]
Put $F=\max_i f_i$.  The affine tails just analyzed show that any negative
endpoint is $o(T_f)$, while
$T_f=(F(r)+F(-r))/2+O(1)$.  Convexity of the $n$-th power and
$a^n+b^n\leq(a+b)^n$ for nonnegative $a,b$ give
\[
 \theta T_f(r)^n+o(T_f^n)
 \leq\frac12\sum_{\delta=\pm1}\sum_i\alpha_i f_i(\delta r)^n
 \leq2^{n-1}\Bigl(\sum_i\alpha_i\Bigr)T_f(r)^n+o(T_f^n).
\]
Combining the last two displays proves both inequalities.
\end{proof}

\begin{leanproof}
\leanfile{secondMain_ordinaryFermat} in \leanfile{Curves/Fermat.lean}.
\end{leanproof}

\begin{example}[5.9 --- Necessity of the growth hypothesis]
\paperlabel{ex:5.9}
\lean{NthTropicalNevanlinna.example59_ambientReciprocalCounting}
\leanok
\uses{thm:5.8}
For $f=[0:2x]$ and $\mathcal P(x_0,x_1)=x_0^2+2x_1^2$,
\[
 T_f(r)=r,\qquad \mathcal P\circ f=8x^2,\qquad
 \sum_{j=1}^2N^{(j)}(r,-\mathcal P\circ f)=8r^2.
\]
Thus $r/T_f(r)=1$, so Theorem 5.8 does not apply.
\end{example}

\begin{proof}
\leanok
The coordinate maximum is $\max(0,2x)$, whose endpoint mean is $r$.  The
function $8x^2$ has only a second-order root at $0$, of multiplicity $16$;
the factor $1/2$ in the integrated count gives $8r^2$.
\end{proof}

\begin{leanproof}
\leanfile{example59_ambientReciprocalCounting} in
\leanfile{Curves/FermatExamples.lean}.
\end{leanproof}

\begin{example}[5.10 --- Sharp constants]
\paperlabel{ex:5.10}
\lean{NthTropicalNevanlinna.example510_concrete_conclusions}
\leanok
\uses{thm:5.8}
Define
\[
 f_1(x)=\begin{cases}0,&x\leq1,\\
 (2k-1)x-(2k^2-1),&2k-1\leq x<2k+1, k\geq1,
 \end{cases}
\]
and
\[
 f_2(x)=\begin{cases}0,&x\leq2,\\
 2kx-2k(k+1),&2k\leq x<2k+2, k\geq1.
 \end{cases}
\]
For $g=[0:f_1]$, $h=[f_1:f_2]$, and $P=x_0+x_1$,
\begin{align*}
 T_g(r)&=N^{(1)}(r,-P\circ g)=\frac14r^2+o(r^2),\\
 T_h(r)&=\frac14r^2+o(r^2),\qquad
 N^{(1)}(r,-P\circ h)=\frac12r^2+o(r^2).
\end{align*}
Thus the lower and upper constants in Theorem 5.8 are sharp.
\end{example}

\begin{proof}
\leanok
Directly on the defining intervals,
$f_1(r)=r^2/2+O(r)$, $f_2(r)=r^2/2+O(r)$, and their maximum has the same
asymptotic; all three functions vanish on the negative tail.  Hence the
Cartan endpoint average for either curve is $r^2/4+o(r^2)$.  The functions
$f_1$ and $f_1+f_2$ are entire, so Jensen identifies their reciprocal counts
with their endpoint means.  These are respectively $r^2/4+o(r^2)$ and
$r^2/2+o(r^2)$, proving both assertions.
\end{proof}

\begin{leanproof}
\leanfile{example510_concrete_conclusions} in
\leanfile{Curves/FermatExamples.lean}.
\end{leanproof}

\section{Casoratian truncation}

\begin{definition}[Shifts, partial counts, and Casoratian]
\paperlabel{def:casoratian-data}
\lean{NthTropicalNevanlinna.forwardShift,NthTropicalNevanlinna.partialIntegratedCounting,NthTropicalNevanlinna.tropicalCasoratian}
\leanok
\uses{def:nevanlinna-quantities}
For $k\geq0$, put $\bar g^{[k]}(x)=g(x+k)$.  For a Borel interval $A$, the
partial integrated count is
\[
 N_A^{(j)}(r,f)=\frac12
 \sum_{\substack{x\in A\cap(-r,r)\\ \omega_f^{(j)}(x)<0}}
 |\omega_f^{(j)}(x)|(r-|x|)^j.
\]
For entire $g_0,\ldots,g_m$, define
\[
 C_0(g_0,\ldots,g_m)(x)=
 \max_{\pi\in S_{m+1}}\sum_{i=0}^m g_i(x+\pi(i)).
\]
\end{definition}

\begin{theorem}[Finite-root first-order identity]
\paperlabel{thm:finite-root-identity}
\lean{NthTropicalNevanlinna.FiniteFirstOrderRootData.integratedRootCounting_eq_affine,NthTropicalNevanlinna.FiniteRootCasoratianData.countingDifference_eventually_constant,NthTropicalNevanlinna.FiniteRootCasoratianData.casoratianCounting_isEquivalent_coordinateCountingSum,NthTropicalNevanlinna.finiteRootCasoratian_countingDifference_isBigO_one,NthTropicalNevanlinna.finiteRootCasoratian_firstOrder_identity,NthTropicalNevanlinna.finiteRootCasoratian_firstOrder_multiplicative}
\leanok
\uses{thm:3.5,def:casoratian-data}
Let $f=[f_0:\cdots:f_m]$ be a first-order tropical holomorphic curve whose
coordinates have finite root sets, possibly empty.  Its Casoratian
$C_0=C_0(f_0,\ldots,f_m)$ is entire of order at most one, has finitely many
roots, and may be constant.  Write the eventual slopes as $k_i^\pm,k_C^\pm$
and put $a=\frac12\sum_i(k_i^+-k_i^-)\geq0$.  Then
\[
 k_C^+-k_C^-=\sum_i(k_i^+-k_i^-),
\]
and there are constants $M_i,M_C$ such that, for all sufficiently large $r$,
\[
 N(r,-f_i)=\frac{k_i^+-k_i^-}{2}r-M_i,\qquad
 N(r,-C_0)=ar-M_C.
\]
Consequently
\[
 N(r,-C_0)-\sum_{i=0}^mN(r,-f_i)=O(1),
\]
and there is a function $\varepsilon(r)\to0$ such that, eventually,
\[
 N(r,-C_0)=\left(\sum_{i=0}^mN(r,-f_i)\right)(1+\varepsilon(r)).
\]
Here the minus signs indicate root counts, using the paper's pole-counting
convention for $N(r,h)$.
\end{theorem}

\begin{proof}
\leanok
Choose a positive bound beyond all coordinate roots; this also makes sense
when some or all root sets are empty.  On either tail, the first-order
multiplicity, hence every change of slope, vanishes.  Continuity gives affine
tails for each coordinate.  Every permutation summand defining $C_0$ has
right slope $\sum_i k_i^+$ and left slope $\sum_i k_i^-$; their finite maximum
has the same tail slopes.  Thus all Casoratian roots lie in a bounded
interval, and local finiteness makes their set finite.

Jensen applied at two radii beyond all roots identifies total root
multiplicity with the difference of the tail slopes.  Put
\[
 M_i=\frac12\sum_{x\in X_i}\omega_{f_i}(x)|x|,\qquad
 M_C=\frac12\sum_{x\in X_C}\omega_{C_0}(x)|x|,
\]
where $X_i,X_C$ are the finite root sets.  Expanding the integrated counts
gives the exact affine formulas.  Their difference is the constant
$\sum_iM_i-M_C$, proving the $O(1)$ estimate.

If $a>0$, the sum of coordinate root counts grows linearly, so the constant
difference divided by that sum tends to zero.  If $a=0$, nonnegativity of all
root multiplicities forces every coordinate and Casoratian root count to
vanish; take $\varepsilon=0$.  No division by a zero count is needed.
\end{proof}

\begin{leanproof}
\leanfile{FiniteFirstOrderRootData.ofFiniteRoots} constructs affine tails from
finite root sets in \leanfile{Truncated/FiniteRoots.lean}.
The curve-facing statements
\leanfile{finiteRootCasoratian_countingDifference_isBigO_one},
\leanfile{finiteRootCasoratian_firstOrder_identity}, and
\leanfile{finiteRootCasoratian_firstOrder_multiplicative} in
\leanfile{Truncated/ShiftCounting.lean} construct the Casoratian data internally.
\end{leanproof}

\begin{lemma}[6.1 --- Casoratian common factor]
\paperlabel{lem:6.1}
\lean{NthTropicalNevanlinna.tropicalCasoratian_commonFactor,NthTropicalNevanlinna.casoratian_commonFactor_of_entire}
\leanok
\uses{def:casoratian-data}
For tropical entire $f_0,\ldots,f_m,g$,
\[
 C_0(f_0\otimesT g,\ldots,f_m\otimesT g)
 =g\otimesT\bar g\otimesT\cdots\otimesT\bar g^{[m]}
  \otimesT C_0(f_0,\ldots,f_m).
\]
\end{lemma}

\begin{proof}
\leanok
Expand the Casoratian as the maximum over permutations.  Every permutation
uses each shift $0,\ldots,m$ exactly once, so the shifted copies of $g$ factor
out as the same additive term in every summand.  Maximum commutes with adding
this common term.
\end{proof}

\begin{leanproof}
\leanfile{tropicalCasoratian_commonFactor} and
\leanfile{casoratian_commonFactor_of_entire} in
\leanfile{Truncated/Casoratian.lean}.
\end{leanproof}

\begin{theorem}[6.2 --- Truncated second main theorem]
\paperlabel{thm:6.2}
\lean{NthTropicalNevanlinna.truncatedSecondMain,NthTropicalNevanlinna.truncatedSecondMainPaperStatement_verified}
\leanok
\uses{thm:3.5,cor:4.8,prop:5.4,lem:6.1}
Let $f=[f_0:\cdots:f_m]$ be a nonconstant reduced $n$-th tropical
holomorphic curve.  Suppose that for some $\ell$ every $f_i-f_\ell$ is
well-defined, that $\rho_2(T_f)=\rho_2<1$, and that
$0<\sigma<1-\rho_2$.  Then, outside a set of finite logarithmic measure,
\begin{align*}
&\left|\sum_{i=0}^m\sum_{j=1}^nN^{(j)}(r,-f_i)
-\sum_{j=1}^nN^{(j)}\!\left(r,-C_0(f_0,\ldots,f_m)\right)\right|\\
&\quad\leq
\sum_{j=1}^n\left|N^{(j)}\!\left(r,C_0(f_0,\ldots,f_m)\right)
-\sum_{i=0}^mN^{(j)}(r,\bar f_i^{[i]})\right|\\
&\qquad+
\sum_{i=0}^m\sum_{j=1}^{\lfloor n/2\rfloor}
\left|N_{[-2i,0]}^{(2j)}(r,-\bar f_i^{[i]})
-N_{[-i,i]}^{(2j)}(r,-f_i)\right|
+o\!\left(\frac{T_f(r)}{r^\sigma}\right).
\end{align*}
\end{theorem}

\begin{proof}
\leanok
For a root $x$ of $f_i$, translation sends it to $x-k$.  The signed
multiplicity is unchanged for odd $j$ and whenever $x$ and $x-k$ have the
same sign.  Thus only even orders crossing the origin contribute a defect,
which is exactly
\[
 N_{[-2k,0]}^{(2j)}(r,-\bar f_i^{[k]})
 -N_{[-k,k]}^{(2j)}(r,-f_i).
\]
Partition the remaining roots into the central region and the four boundary
strips determined by $r-k,r,r+k$.  Monotonicity of root counting sandwiches
their total contribution between the corresponding counts at $r-k$ and
$r+k$.  Proposition 5.4 dominates every coordinate quotient characteristic
by $T_f+O(1)$; Lemma 4.6 then makes each fixed-radius increment
$o(T_f/r^\sigma)$.  Summing the finitely many $i,j,k$ proves the shifted-count
identity used in the displayed right side.

By Lemma 6.1, subtracting the reference coordinate $f_\ell$ from all
coordinates only adds the common shifted factor
$\sum_{k=0}^m f_\ell(x+k)$ to the Casoratian.  At either endpoint $\delta r$,
choose a maximizing permutation $\pi_\delta$.  Theorem 4.7, applied to each
well-defined quotient $f_i-f_\ell$, gives
\[
 \left|C_0(f_0,\ldots,f_m)(\delta r)
       -\sum_{i=0}^m\bar f_i^{[i]}(\delta r)\right|
 =o(T_f(r)/r^\sigma).
\]
Indeed, after removing the common factor the difference is bounded by the
finite sum of shift differences between $\pi_\delta(i)$ and $i$.

Finally apply the $n$-th Jensen identity to $C_0$ and separately to every
$\bar f_i^{[i]}$.  Insert the shifted-count identity from the first paragraph
and the endpoint estimate from the second, then apply the triangle inequality.
This is exactly the asserted inequality.  Nonconstancy, convexity, and
Proposition 5.4 imply the characteristic scale tends to infinity and absorb
the finitely many $O(1)$ terms; no additional growth assumption is used.
\end{proof}

\begin{leanproof}
\leanfile{truncatedSecondMain} and
\leanfile{truncatedSecondMainPaperStatement_verified} in
\leanfile{Truncated/Main.lean}; shift transport and the exact even correction are
in \leanfile{Truncated/ShiftCounting.lean}.
\end{leanproof}

\begin{example}[Higher-order correction terms are necessary]
\paperlabel{ex:section6-obstruction}
\lean{NthTropicalNevanlinna.section6Counterexample_exact_obstruction}
\leanok
\uses{thm:6.2}
Let
\[
 f_0(x)=\begin{cases}-x^2,&x\leq0,\\x^2,&x>0,\end{cases}
 \qquad
 f_1(x)=\begin{cases}x^2,&x\leq0,\\-x^2,&x>0.\end{cases}
\]
Then $f=[f_0:f_1]$ is a well-defined second-order curve and
\[
 C_0(f_0,f_1)(x)=\begin{cases}
 -2x-1,&x\leq-1,\\
 2x^2+2x+1,&-1<x\leq0,\\
 2x+1,&x>0.
 \end{cases}
\]
Moreover $T_f(r)=r^2$, the coordinate reciprocal counts vanish, while
\[
                 \sum_{j=1}^2N^{(j)}(r,-C_0)=r^2.
\]
Hence the two explicit error sums in Theorem 6.2 cannot in general be
replaced by $o(T_f/r^\sigma)$ when $n\geq2$.
\end{example}

\begin{proof}
\leanok
The coordinate maximum is $x^2$, giving $T_f(r)=r^2$.  The negatives of both
coordinates are second-order entire, so their reciprocal counts are zero.
Taking the maximum of the two Casoratian permutation sums yields the displayed
three pieces.  Its only reciprocal-count contribution is the normalized
second-order root at $0$; direct substitution gives integrated count $r^2$.
\end{proof}

\begin{leanproof}
\leanfile{section6Counterexample_exact_obstruction} in
\leanfile{Counterexamples.lean}.
\end{leanproof}

\begin{corollary}[6.3 --- First-order truncated relation]
\paperlabel{cor:6.3}
\lean{NthTropicalNevanlinna.firstOrderTruncatedSecondMainPaperStatement_verified}
\leanok
\uses{thm:6.2}
For a nonconstant reduced first-order curve satisfying the same hyper-order
condition,
\[
 \sum_{i=0}^mN(r,-f_i)
 =N\!\left(r,-C_0(f_0,\ldots,f_m)\right)
 +o\!\left(\frac{T_f(r)}{r^\sigma}\right)
\]
outside a set of finite logarithmic measure.
\end{corollary}

\begin{proof}
\leanok
At first order there is no even-order correction.  The shifted coordinate
functions and the Casoratian are tropical entire, so the two explicit pole
count errors in Theorem 6.2 vanish.  The remaining term is the stated
little-$o$ remainder.
\end{proof}

\begin{leanproof}
\leanfile{firstOrderTruncatedSecondMainPaperStatement_verified} in
\leanfile{Truncated/Main.lean}.
\end{leanproof}

\begin{example}[The hyper-order condition in Corollary 6.3 is necessary]
\paperlabel{ex:cor6.3-growth}
\lean{NthTropicalNevanlinna.corollary63_CasoratianReciprocalCounting_isEquivalent_scaledCharacteristic,NthTropicalNevanlinna.corollary63_growthHypothesis_counterexample}
\leanok
\uses{prop:3.8,cor:6.3}
Let $f=[0:e_\alpha]$ with $\alpha>1$.  Then
$T_f(r)=T(r,e_\alpha)+O(1)$ has hyper-order $1$,
$e_\alpha(x+1)=\alpha e_\alpha(x)$, and
\[
 C_0(0,e_\alpha)=\alpha e_\alpha.
\]
Consequently
\begin{align*}
 N(r,-0)+N(r,-e_\alpha)&\sim T(r,e_\alpha),\\
 N(r,-C_0(0,e_\alpha))&\sim\alpha T(r,e_\alpha),
\end{align*}
so the conclusion of Corollary 6.3 fails.
\end{example}

\begin{proof}
\leanok
The shift equation follows directly from the floor-function definition of
$e_\alpha$.  Of the two Casoratian permutation sums, nonnegativity makes the
shifted hyper-exponential the maximum, hence the Casoratian equals
$\alpha e_\alpha$.  Proposition 3.8 gives hyper-order $1$.  Jensen gives the
first asymptotic, and positive scalar homogeneity of multiplicities multiplies
the second count by $\alpha$.  Since $\alpha>1$, the two leading terms cannot
differ by the little-$o$ remainder required in Corollary 6.3.
\end{proof}

\begin{leanproof}
\leanfile{corollary63_CasoratianReciprocalCounting_isEquivalent_scaledCharacteristic}
and \leanfile{corollary63_growthHypothesis_counterexample} in
\leanfile{Counterexamples.lean}.
\end{leanproof}
